\documentclass[12pt,a4paper]{article}
\usepackage[margin=25mm, headheight=15pt, headsep=10pt]{geometry}
\usepackage{fontspec}
\usepackage[table]{xcolor} % Note: table option is required for rowcolors
\usepackage{titlesec}
\usepackage{tocloft}
\usepackage{fancyhdr}
\usepackage{booktabs}
\usepackage{tcolorbox}
\tcbuselibrary{skins, breakable}
\usepackage[colorlinks=true, linkcolor=emerald, urlcolor=emerald, bookmarks=true]{hyperref}

\definecolor{emerald}{RGB}{0,102,51}
\definecolor{darkred}{RGB}{139,0,0}
\definecolor{lightgray}{RGB}{245,245,245}

\usepackage[nil]{babel}
\babelprovide[import, main]{bengali}
\babelprovide[import]{arabic}
\babelfont{rm}[Renderer=HarfBuzz,Scale=1.0]{Noto Serif Bengali}
\babelfont{sf}[Renderer=HarfBuzz]{Noto Sans Bengali}
\babelfont{tt}[Renderer=HarfBuzz]{Noto Sans Bengali}
\babelfont[arabic]{rm}[Renderer=HarfBuzz,Scale=1.1]{Noto Naskh Arabic}

% Section styling
\titleformat{\section}{\Large\bfseries\color{emerald}}{\thesection.}{0.5em}{}
\titleformat{\subsection}{\large\bfseries\color{emerald!85!black}}{\thesubsection.}{0.5em}{}
\titleformat{\subsubsection}{\normalsize\bfseries\color{emerald!70!black}}{\thesubsubsection.}{0.5em}{}
\titlespacing*{\section}{0pt}{18pt}{8pt}

% Fancy Header/Footer setup
\pagestyle{fancy}
\fancyhf{} % clear all header and footer fields
\fancyhead[L]{\color{emerald}\small\textbf{\nouppercase{\leftmark}}}
\fancyfoot[L]{\scriptsize\color{black!50}{\lat{AI}}-সংকলিত, আলেম-যাচাইকৃত নয়}
\fancyfoot[C]{\color{emerald!70!black}\thepage}
\renewcommand{\headrulewidth}{0.4pt}
\renewcommand{\headrule}{\hbox to\headwidth{\color{emerald}\leaders\hrule height \headrulewidth\hfill}}
\renewcommand{\footrulewidth}{0pt}

% Custom boxes
% 1. Ayat/Hadith Box
\newtcolorbox{ayatbox}[1][]{
    enhanced, breakable, 
    colback=emerald!5, 
    colframe=emerald, 
    boxrule=0.5pt, 
    arc=3pt, 
    left=10pt, right=10pt, top=10pt, bottom=10pt,
    #1
}

% 2. Quote/Translation Box
\newtcolorbox{quotebox}[1][]{
    enhanced, breakable,
    frame hidden,
    colback=lightgray,
    borderline west={3pt}{0pt}{emerald},
    left=15pt, right=10pt, top=10pt, bottom=10pt,
    sharp corners,
    #1
}

% 3. AI-generated notice box
\newtcolorbox{warnbox}[1][]{
    enhanced, breakable,
    colback=red!4,
    colframe=darkred,
    boxrule=0.8pt,
    arc=3pt,
    left=10pt, right=10pt, top=10pt, bottom=10pt,
    #1
}

% Commands
\newcommand{\arab}[1]{\foreignlanguage{arabic}{#1}}
\newcommand{\kib}[1]{\textcolor{emerald}{\textbf{#1}}}

% Latin (English) fallback: Noto Serif Bengali has NO Latin glyphs
\newfontfamily\latfont[Renderer=HarfBuzz]{Noto Serif}
\newcommand{\lat}[1]{{\latfont #1}}

\setlength{\parskip}{5pt}
\setlength{\parindent}{0pt}

% Fix for missing bullet in Noto Serif Bengali
\renewcommand{\labelitemi}{$\bullet$}



\begin{document}
\begin{center}{\Large\bfseries\color{emerald} পার্ট ৭খ — নামাজের ভুল ও প্রতিকার: ভুল → বিধান রেফারেন্স}\end{center}
\vspace{1em}
\section*{পার্ট ৭খ — নামাজের ভুল ও প্রতিকার: ভুল $\rightarrow$ বিধান রেফারেন্স}

\begin{warnbox}
 \textbf{গুরুত্বপূর্ণ নোটিশ (\lat{Important} \lat{Notice}):} এই কন্টেন্টটি \lat{AI} (কৃত্রিম বুদ্ধিমত্তা) দিয়ে প্রস্তুত ও সংকলিত। \lat{AI} কঠোর সূত্র-মান নীতি মেনে শুধু নির্ভরযোগ্য উৎস থেকে তথ্য সংগ্রহ করেছে — কেবল সহীহ/হাসান হাদিস ও সহীহ সনদের আসার রাখা হয়েছে, দুর্বল সনদ বাদ দেওয়া হয়েছে। তবে এটি কোনো আলেম বা মুহাদ্দিস কর্তৃক যাচাইকৃত নয়।
\end{warnbox}

\begin{quote}
\textbf{সম্পর্কিত ফাইল:} পার্ট ৩ (শ্রেণিবিভাগ), পার্ট ৫খ (সাহু সিজদা), পার্ট ৭ক (ভঙ্গকারী)।
\end{quote}

\begin{quote}
\textbf{যে প্রশ্নের উত্তর এই অংশে:} নামাজে ভুল হলে \textbf{ধাপে ধাপে} কী করবেন? রাকাত কম-বেশি, রুকু-সিজদা বাদ, মাসবুকের ভুল — সব কেস-স্টাডি।
\end{quote}

\medskip\hrule\medskip

\subsection*{১. ভুল-প্রতিকারের মূল সূত্র (নীতিমালা)}

\begin{enumerate}
\item \textbf{রুকন (ফরজ) ছুটে গেলে:} মনে পড়া মাত্রই — যদি \textbf{পরের রুকনে যাওয়ার আগে} মনে পড়ে, ফিরে গিয়ে পূরণ করুন (নামাজ চলবে)। যদি \textbf{পরের রুকনে পৌঁছে} যান, পুরো রাকাত পুনরায় + শেষে সাহু সিজদা।
\item \textbf{ওয়াজিব (হানাফি) ছুটে গেলে:} সাহু সিজদা — সালামের আগে বা পরে (পার্ট ৫খ)।
\item \textbf{সন্দেহ হলে:} নিশ্চিত ন্যূনতম ধরে + সাহু সিজদা।
\end{enumerate}
\medskip\hrule\medskip

\subsection*{২. কেস-স্টাডি টেবিল (ভুল $\rightarrow$ সমাধান)}

\begin{table}[h]\centering\small
\begin{tabular}{lll}
\# & ভুল & সমাধান (প্রধানত হানাফি-জমহুরের মিশ্র) \\
\hline
১ & ভুলে রুকু বাদ গেল (সিজদায় ঢুকে পড়েছেন) & পরের রুকনে পৌঁছে গেলে: ১ রাকাত বাতিল — পূর্ণ রাকাত ফিরে, শেষে সাহু সিজদা \\
২ & সিজদা বাদ গেল (পরের সিজদা-রাকাতে) & একই নিয়ম: ফিরে গিয়ে বাদ-পড়া সিজদা (পরের রাকাতে) + শেষে সাহু সিজদা \\
৩ & ফাতিহা/সূরা ভুলে ছুটল (ওয়াজিব) & সাহু সিজদা \\
৪ & তাশাহহুদ ভুলে (প্রথম বৈঠক) ছাড়ল & সাহু সিজদা (হানাফি); জমহুরে ফরজ বাদ — নামাজ বাতিল (বিতর্কিত) \\
৫ & রাকাত ৪-এর বদলে ৩ পড়ল (নিশ্চিত) & ৪র্থ রাকাত যোগ + সাহু সিজদা \\
৬ & ৪-এর বদলে ৫ পড়ল (নিশ্চিত) & সাহু সিজদা দিয়ে সংশোধন (নামাজ বাতিল নয় — বুখারি ৪০১) \\
৭ & সন্দেহ: ৩ না ৪? & ৩ ধরে নিন + সাহু সিজদা (নিশ্চিত-ন্যূনতম) \\
৮ & সালামের পর মনে পড়ল — রাকাত কম & এক রাকাত যোগ $\rightarrow$ তাশাহহুদ $\rightarrow$ সালাম $\rightarrow$ সাহু সিজদা (সালাম-পরে-যদি ইত্যাদি মাযহাবগত) \\
৯ & ভুলে ১ রাকাত কম পড়ে সালাম দিলেন & ইমাম-নিয়মে: সোজা উঠে বাকি রাকাত $\rightarrow$ শেষ তাশাহহুদ $\rightarrow$ সালাম $\rightarrow$ সাহু সিজদা \\
১০ & মাসবুক — ইমামের ২য় রাকাতে ঢুকল & ইমামের সাথে তাশাহহুদ $\rightarrow$ ইমাম সালামের পর ২ রাকাত পূর্ণ $\rightarrow$ নিজের ভুলের জন্য সাহু সিজদা (নিজের রাকাতে) \\
\hline
\end{tabular}
\end{table}

\begin{quote}
 \textbf{গুরুত্বপূর্ণ সতর্কতা:} উপরের টেবিল \textbf{সাধারণ রূপরেখা} — প্রতিটি কেসে মাযহাবগত সূক্ষ্মতা আছে (যেমন হানাফিতে তাশাহহুদ-প্রথম বৈঠক = ওয়াজিব-সাহু; জমহুরে ফরজ-বাতিল)। \textbf{নিশ্চিত জ্ঞান ছাড়া নিজের ভুল নিজে মীমাংসা না করা} — ইমাম/আলেমের কাছে জিজ্ঞেস করা সর্বোত্তম।
\end{quote}

\medskip\hrule\medskip

\subsection*{৩. মাসবুকের জটিল কেস}

\subsubsection*{কেস এ: মাসবুক ২ রাকাত পেয়েছে — নিজের ভুলের সাহু সিজদা}

\begin{itemize}
\item ইমামের পেছনে ২ রাকাত (যোহর) $\rightarrow$ ইমাম সালাম $\rightarrow$ মাসবুক উঠে ২ রাকাত নিজে $\rightarrow$ নিজের ২ রাকাতের শেষে নিজের ভুল থাকলে \textbf{নিজের সালামের আগে সাহু সিজদা}। (ইমামের পেছনের ২ রাকাতের ভুল — ইমামের সাহু সিজদাতেই কভার)
\end{itemize}
\subsubsection*{কেস বি: মাসবুক ৩য় রাকাতে ঢুকল — হানাফি-জমহুরের ভিন্নতা}

\begin{itemize}
\item ইমামের সাথে ৩য় রাকাত (যা তার নিজের ২য়?) — এখানে জটিলতা; \textbf{সাধারণ নিয়ম:} যা ইমামের সাথে পেল = তার শেষ রাকাত; বাকি নিজে পূর্ণ। মাঝের রাকাতে তাশাহহুদ বসার সিদ্ধান্ত মাযহাবভেদে ভিন্ন। — \textbf{পরামর্শ:} জটিল অবস্থায় আলেমের কাছে কেস-ভিত্তিক প্রশ্ন।
\end{itemize}
\medskip\hrule\medskip

\subsection*{৪. নামাজের ভুল-সংশোধনের "সোনালি নিয়ম"}

\begin{enumerate}
\item \textbf{নিশ্চিত ভুল} $\rightarrow$ মীমাংসা $\rightarrow$ সাহু সিজদা (প্রয়োজনে)।
\item \textbf{সন্দেহ} $\rightarrow$ নিশ্চিত ন্যূনতম $\rightarrow$ সাহু সিজদা।
\item \textbf{ভুলে কথা/খাওয়া} $\rightarrow$ ভুল হলে মাযহাবগত; নিয়ম: নামাজ বাতিলের আশঙ্কায় \textbf{যথাসম্ভব ইমামের সাথে} — শেষে ইমামকে/আলেমকে বলা।
\item \textbf{নিজের ভুল আত্ম-সমাধানে চিন্তিত} — হতাশ না হয়ে \textbf{নতুন নামাজ} ও \textbf{সাহু সিজদা-জ্ঞান} দিয়ে সামনে এগোন।
\end{enumerate}
\begin{quote}
\textbf{মূল শিক্ষা:} নামাজের ভুল মানুষমাত্রই হয় — নবীজি (সা.)-ও (মানুষ হিসেবে) ভুল করেছেন এবং সাহু সিজদা শিখিয়েছেন। \textbf{ভুলে ভয় নয়; শেখা ও সংশোধনই} — এটাই এই অংশের লক্ষ্য।
\end{quote}

\medskip\hrule\medskip

\subsection*{৫. রেফারেন্স}

\begin{itemize}
\item সহীহ বুখারি — ৪০১, ৪০৪, ৬৩৫
\item সহীহ মুসলিম — ৫৭২, ৫৭৩, ৫৭৪
\item সুনানে আবু দাউদ — ১০১৯-১০২০, ১০২৪
\item সুনানে ইবনে মাজাহ — ১২১১-১২১২
\item আল-ফিকহুল ইসলামি ওয়া আদিল্লাতুহু (আল-যুহাইলি) — কেস-ভিত্তিক মাযহাব-সিদ্ধান্ত
\end{itemize}
\begin{quote}
\textbf{দ্রষ্টব্য:} এই অংশের টেবিল \textbf{সাধারণ দিকনির্দেশনা} — ব্যক্তিগত কেসে নির্ভরযোগ্য আলেমের ফতোয়া সর্বোত্তম।
\end{quote}

\end{document}
